\textbf{IMPLEMENTATION}

This chapter describes how the five architectural components introduced
earlier were realised in practice: the repository and
development-environment organisation, the dataset and baseline
pipelines, the QLoRA fine-tuning procedure, the BMAD agentic correction
loop and its underlying inference mechanism, and the evaluation tooling
used to produce the results reported in Chapter 5. Particular attention
is given to Section 4.5, which documents both the original
implementation of the agentic loop's inference layer and a subsequent
architectural revision undertaken after the original design proved
unreliable under sustained, unattended execution --- a revision that is
reported here in full because the diagnosis and correction of that
unreliability is itself a substantive piece of engineering work
underpinning the validity of the results in Chapter 5.

\textbf{4.1 Development Environment and Repository Organisation}

Development was distributed across two machines according to
computational demand. Mac A, an Apple M4 Pro with 48 GB of unified
memory, hosted all locally executed model inference: QLoRA fine-tuning,
the BMAD agentic loop, and the statistical evaluation. Mac B, an
Intel-based machine with 16 GB of RAM and no CUDA-capable GPU, was
reserved for work that depends only on network calls to external APIs
rather than local model execution --- principally dataset generation and
the two zero-shot commercial baselines. This division reflects a
deliberate architectural decision established early in the project and
confirmed empirically during development: Mac B's hardware is unsuited
to hosting a 3.8--4-billion-parameter transformer for repeated,
low-latency inference, whereas its Intel CPU is entirely adequate for
issuing HTTP requests to a commercial LLM API and storing the response.

The codebase is organised as a single Git repository with directories
corresponding to each pipeline stage --- data/, baselines/,
fine\_tuning/, agentic\_loop/, evaluation/ and results/ --- and is
version-controlled across four branches: main, which holds the dataset,
baseline, and fine-tuning code common to all subsequent work;
phase6-eval, an early implementation of the agentic loop later found to
have a systemic reliability defect (Section 4.5.1) and retained for
historical record rather than active use; phase6-rebuild, the corrected
implementation used to produce the results in Chapter 5; and
final-submission, into which phase6-rebuild has been merged as the basis
for this dissertation.

\textbf{4.2 Dataset Generation and Preparation}

The 279-story dataset was produced by a generation script that
synthesises Jira-style user stories paired with a Cypress and a
Playwright reference script for each story, balanced across sixteen
functional categories (authentication, CRUD operations, form validation,
responsive design, and others) and three complexity tiers (simple,
medium, complex). A separate splitting utility partitions the unified
dataset into framework-isolated JSON Lines files, one containing only
Cypress reference scripts and one containing only Playwright reference
scripts, so that fine-tuning and baseline evaluation for a given
framework never has visibility of the other framework's syntax during
training or prompting.

\textbf{4.3 Zero-Shot Baseline Inference}

The baseline runner issues one API request per user story to each of the
two commercial models under evaluation, GPT-4o-mini and Claude Haiku,
and persists the response together with prompt and completion token
counts, latency, and a timestamp as an individual JSON record. A
per-model, per-framework summary file records aggregate completion
counts, allowing an interrupted baseline run to resume from the last
completed record rather than reissuing already-completed API calls. A
third candidate baseline, Gemini Flash, was evaluated for inclusion but
excluded after its free-tier request quota proved insufficient for the
558 calls required across both frameworks; this decision is recorded in
the codebase rather than silently omitted.

\textbf{4.4 QLoRA Fine-Tuning}

Each of the four adapters --- Phi-3 Mini and Gemma 4 E4B, one instance
per framework --- was produced by a dedicated fine-tuning script
applying Low-Rank Adaptation {[}6{]} under 4-bit quantisation (QLoRA)
{[}7{]} to the frozen base model, using the framework-isolated training
split described in Section 4.2. Training and inference both execute on
Apple Silicon through PyTorch's Metal Performance Shaders (MPS) backend,
which required three specific adaptations relative to a conventional
CUDA training configuration, identified during the Phase 4 diagnostic
experiments reported at the mid-semester stage: the bfloat16 numeric
format {[}20{]} must be used in place of float16, since PyTorch's
gradient scaler silently skips optimiser steps under float16 on MPS;
HuggingFace TRL's {[}21{]} completion-only masking collator produced NaN
losses under certain batch configurations and required adjustment; and
Python's standard output buffering must be disabled at the interpreter
level to obtain real-time loss visibility during training. These
adaptations, once established, reduced the iteration time for subsequent
fine-tuning experiments materially, as reflected in the substantially
shorter Phi-3 Mini training times reported in the mid-semester results
(under one hour per framework, against 3.78 and 5.71 hours for Gemma 4
E4B).

\textbf{4.5 The BMAD Agentic Correction Loop}

The agentic loop is the central engineering artefact of this
dissertation. It implements a Build--Measure--Assess--Decide cycle: for
a given user story, the Build step generates a candidate script from the
appropriate fine-tuned adapter; the Measure step scores the candidate
against a composite quality rubric; the Assess step compares that score
to an acceptance threshold; and, if the threshold is not met, the Decide
step constructs targeted corrective feedback and returns control to
Build for a further attempt, up to a configurable maximum of three
iterations, after which the best-scoring candidate across all attempts
is retained regardless of whether it cleared the threshold.

\textbf{4.5.1 Original Implementation and Its Failure Modes}

The loop was first implemented as two separate processes communicating
over HTTP: a FastAPI/uvicorn inference server that loaded a fine-tuned
adapter and exposed a single generation endpoint on a local port, and a
client process implementing the Build--Measure--Assess--Decide logic
that issued one HTTP request per generation attempt. This separation is
a conventional and reasonable design for a service intended to handle
concurrent external callers; however, subjecting it to sustained,
unattended execution across the full 279-story dataset for all four
adapters --- a workload spanning tens of continuous hours per adapter
--- exposed four distinct reliability defects, each diagnosed and
corrected in turn.

\begin{itemize}
\item
  A port-binding conflict: a previous inference-server process that had
  not been fully terminated continued to hold the local port, so a fresh
  server process for the next model/framework combination failed to bind
  while the orchestration script's health check nonetheless reported the
  (stale) server as healthy, causing a subsequent generation combination
  to run against the wrong model without failing loudly.
\item
  Read-timeout failures under memory pressure: when overall system
  memory pressure was high, individual generation calls occasionally
  exceeded the client's request timeout even though the server process
  was still alive, and this exception type had not been included in the
  client's retry-on-failure logic, so a single slow response terminated
  the entire multi-hour run.
\item
  An incorrect exception reference in the retry-handling code itself:
  the client's connection-retry logic referenced an exception class that
  does not exist in the networking library in use, so on the one
  occasion a genuine dropped connection occurred, the retry mechanism
  intended to handle it raised a second, unrelated error and crashed
  rather than retrying.
\item
  A defect introduced while fixing a memory-growth issue: an interim
  mitigation for gradual memory growth in the long-lived server process
  (Section 4.5.2) itself contained a variable-scoping error that caused
  every subsequent request to fail; this was caught by validating the
  fix against a live request before redeploying it, rather than relying
  on a syntax check alone, and is recorded here as an example of why
  runtime verification, not just static verification, was adopted as
  standard practice for the remainder of the implementation.
\end{itemize}

Each of these defects shares a common root cause: the HTTP boundary
between the correction loop and the model introduced a class of failure
--- connection state, timeouts, port lifecycle, serialisation of retry
logic --- that has nothing to do with the correctness of the generated
test scripts, yet was capable of terminating a multi-hour unattended
run. This observation motivated a architectural revision rather than a
further round of point fixes.

\textbf{4.5.2 Rearchitected Single-Process Implementation}

The corrected implementation, on which all results in Chapter 5 are
based, removes the HTTP boundary entirely. Model loading, prompt
construction, and generation are performed by a single in-process module
that the Build--Measure--Assess--Decide loop calls directly as a
function, with no server, no port, and no network client involved. This
eliminates the entire class of failure described in Section 4.5.1 by
construction, since there is no longer a connection to drop, a port to
conflict over, or a timeout to elapse. Two further corrections, both
genuine model-behaviour fixes rather than infrastructure fixes, were
carried into the rearchitected implementation:

\begin{itemize}
\item
  Gemma 4 E4B's generation calls previously specified only the
  tokenizer's default end-of-sequence token as a stopping criterion;
  because Gemma 4's chat template delimits a turn with a distinct
  \textless end\_of\_turn\textgreater{} marker, generation would
  occasionally continue past the intended script into a hallucinated
  continuation of the conversation. The corrected implementation passes
  the model-specific turn-boundary token explicitly as an additional
  stopping criterion, and applies a truncate-at-first-marker safety net
  on the decoded output as a second line of defence should the stopping
  criterion still be missed.
\item
  Memory usage on the Metal Performance Shaders backend was found to
  grow gradually across hundreds of sequential generation calls within a
  single long-lived process, eventually contributing to the read-timeout
  failures described above. An explicit garbage-collection pass together
  with the MPS-specific cache-clearing call is now issued after every
  generation, which measurably stabilised memory usage across multi-hour
  continuous runs in subsequent testing.
\end{itemize}

The loop additionally implements a token-budget escalation policy: if a
candidate script is rejected specifically because its braces or
parentheses are unbalanced --- the signature of truncation against the
generation token ceiling rather than a content-quality failure --- the
next correction attempt is issued with a higher token ceiling rather
than an identical one, since re-requesting ``the complete script'' at an
unchanged budget simply reproduces the same cutoff. The ceiling is
capped to bound worst-case latency per record.

\textbf{4.5.3 Execution Reliability Engineering}

Beyond the single-process architectural change, three operational
measures were adopted to make the full-dataset run practical on a single
development machine shared with other applications, rather than a
dedicated server.

\begin{itemize}
\item
  Checkpoint and resume: before generating a script for a given story,
  the loop checks whether a result file for that story identifier
  already exists in the output directory and, if so, skips regeneration.
  This allows a run to be safely interrupted and restarted from the
  point of interruption without discarding completed work, and underlies
  the chunked execution strategy below.
\item
  Chunked execution: rather than attempting all 279 records per adapter
  in a single continuous process lifetime, execution was performed in
  fixed increments of 35 records per model/framework combination, with
  the process exiting and being relaunched between increments. Combined
  with checkpoint/resume, this bounds the maximum continuous process
  lifetime and correspondingly bounds any residual memory growth, while
  the retry-with-restart wrapper in the orchestration script
  additionally re-attempts a failed combination up to three times,
  restarting the process and resuming from checkpoint, before surfacing
  an error rather than continuing silently.
\item
  Sleep prevention: because each chunk could run for several hours
  unattended, the operating system's default idle-sleep behaviour ---
  which would otherwise suspend all running processes after a short
  period of inactivity --- was suppressed for the duration of each chunk
  using a process-scoped sleep-prevention utility tied to the lifetime
  of the pipeline process, so that sleep prevention is released
  automatically when the chunk completes rather than requiring manual
  cleanup.
\end{itemize}

\textbf{4.6 Evaluation Tooling}

Two categories of evaluation tooling were implemented to produce the
results in Chapter 5. Independent syntax validation runs Node.js's
built-in JavaScript/TypeScript parser (node -\/-check) against every
generated script from all eight systems, applying the same stop-marker
truncation used defensively during generation (Section 4.5.2) so that
any residual leaked chat-template text does not distort the validity
measurement, and a corresponding Markdown-code-fence stripping step for
the baseline systems after the fence-related measurement artefact
described in Section 5.2 was identified. Statistical analysis tooling,
built on the SciPy library, loads the per-record scores produced by both
the BMAD loop and the baseline evaluator, joins them by story identifier
to construct the paired dataset described in Section 5.1, and computes
the one-way ANOVA, paired Wilcoxon signed-rank tests, chi-square test of
independence, and category-level error matrices reported throughout
Chapter 5.

\textbf{4.7 Summary}

The implementation delivers all five architectural components introduced
earlier: a framework-isolated dataset pipeline, a resumable zero-shot
baseline evaluator, an MPS-adapted QLoRA fine-tuning procedure, a
single-process BMAD agentic correction loop, and a statistical
evaluation toolchain. The most significant engineering finding to emerge
during implementation was that the original two-process, HTTP-mediated
design for the agentic loop was reliable for short runs but not for the
sustained, multi-hour unattended execution required to process the
complete dataset, and that removing the HTTP boundary entirely ---
rather than continuing to patch individual failure modes within it ---
was the more effective correction. All quantitative results reported in
Chapter 5 were produced using the corrected, single-process
implementation described in Section 4.5.2, executed to completion across
all four model/framework combinations with zero unrecovered failures
once the chunked execution strategy of Section 4.5.3 was adopted.
